\paragraph{Observações.}
\begin{itemize}
  \item \textbf{(ii) $\approx$ (iii):} o acoplamento hidromecânico com ar passivo muda o FS mínimo em $+0{,}001$ a
    $+0{,}03$, exceto no f50 com 90~mm/h ($+0{,}074$: 0,906 contra 0,832), onde a (iii) tem a queda brusca nas duas
    últimas horas da chuva e a (ii) não. O $t_1$ quase não muda (10,9 contra 10,7~h no f10; 20,5 contra 19,3~h no f50).
    O deslocamento máximo do esqueleto elástico é de 1 a 7~cm.
  \item \textbf{(i), ar ativo:} nas chuvas fracas e moderadas do f50 e do f100 o FS cai muito mais: $-0{,}12$ (f50,
    9~mm/h), $-0{,}12$ (f50, 36~mm/h, que passa a romper: 0,919, $t_1 = 15{,}0$~h), $-0{,}16$ (f100, 3,6~mm/h),
    $-0{,}28$ (f100, 9~mm/h) e $-0{,}25$ (f100, 36~mm/h). No f10, com $I_r \le 0{,}25k_s$, quase nada muda (o ar
    escapa: $p_a \le 23$~kPa).
  \item \textbf{(i) nas chuvas fortes:} o FS mínimo fica \emph{acima} da (iii) (f10 com 180~mm/h: 0,896 contra 0,842;
    f50 com 90~mm/h: 0,910 contra 0,832), mas o talude rompe muito antes ($t_1$ = 5,8 contra 10,7~h e 11,3 contra
    19,3~h). Pelo FS mínimo o ar ativo pareceria favorável; pela primeira passagem ele é o mais perigoso. O ar
    comprimido sob a frente reduz a resistência cedo, e depois a menor infiltração (Seção~\ref{sec:trans}) limita a
    queda final.
\end{itemize}
